% Rédaction : BANDE Hamidou ; relecture : KOUMTOUBRE Davy
\chapter{Contexte et exigences}

\section{Le sujet}
Le thème 6 du projet \emph{Développement à base de composants et services Web} demande une
API REST de transactions bancaires construite avec Express.js, alimentée par un workflow KNIME et
supervisée par Prometheus et Grafana. L'enjeu n'est pas une seule technologie mais leur
assemblage : chaque composant a un rôle distinct et doit pouvoir être remplacé ou testé seul.

Notre chaîne fait cinq choses :
\begin{enumerate}
  \item \textbf{générer} 1010 transactions bancaires fictives, dont des erreurs injectées exprès ;
  \item \textbf{contrôler et transformer} ces données avec KNIME : 946 lignes valides, 64 rejetées avec leur motif ;
  \item \textbf{charger} le résultat dans PostgreSQL, sans jamais créer de doublon ;
  \item \textbf{exposer} les transactions par une API REST filtrable, paginée et modifiable ;
  \item \textbf{mesurer} les requêtes, les erreurs et les temps de réponse de l'API.
\end{enumerate}

\section{Exigences de l'exercice 6}
\begin{table}[H]
\centering\small
\begin{tabularx}{\linewidth}{@{}lXl@{}}
\toprule
\textbf{N°} & \textbf{Travail demandé (énoncé)} & \textbf{Chapitre} \\
\midrule
E1 & Générer des transactions bancaires fictives & \ref{chap:donnees} \\
E2 & Construire un workflow KNIME pour contrôler et transformer les données & \ref{chap:knime} \\
E3 & Développer l'API REST et ses filtres & \ref{chap:api} \\
E4 & Exposer les métriques de requêtes, d'erreurs et de temps de réponse & \ref{chap:supervision} \\
E5 & Configurer Prometheus et Grafana & \ref{chap:supervision} \\
E6 & Utiliser Cypress si une interface web est fournie ; sinon, un outil adapté aux tests d'API & \ref{chap:tests} \\
\bottomrule
\end{tabularx}
\caption{Exigences de l'exercice 6 et chapitre qui y répond}
\end{table}

\section{Exigences communes à tous les groupes}
\begin{table}[H]
\centering\small
\begin{tabularx}{\linewidth}{@{}lX@{}}
\toprule
\textbf{N°} & \textbf{Exigence} \\
\midrule
G1 & Dépôt Git public : code et dépendances, scripts de création des bases, configurations, workflows et contrats d'API, README (architecture, prérequis, installation), tests et résultats attendus, présentation PowerPoint. \\
G2 & Installation reproductible sur une machine neuve ; Docker Compose pour les services compatibles ; procédure distincte pour les outils à installation particulière (KNIME). \\
G3 & Au minimum un scénario fonctionnel réussi et un test de bout en bout de la chaîne principale. \\
\bottomrule
\end{tabularx}
\caption{Exigences communes}
\end{table}

\section{Hypothèses retenues}
L'énoncé laisse plusieurs points ouverts. Nos choix, justifiés au chapitre~\ref{chap:architecture} :
\begin{itemize}
  \item la base n'est pas imposée : nous utilisons \textbf{PostgreSQL 16} ;
  \item « l'API et ses filtres » : consultation filtrée \emph{et} création, modification, suppression ;
  \item KNIME écrit deux fichiers CSV, vérifiés puis chargés par un chargeur idempotent, plutôt que d'écrire directement dans la base ;
  \item aucune interface web n'est demandée : les tests utilisent Jest, Supertest, Postman (Newman) et un test de bout en bout, pas Cypress ;
  \item l'authentification est hors périmètre (piste d'amélioration au chapitre~\ref{chap:limites}).
\end{itemize}
